\documentclass[sigconf,nonacm]{acmart}

\usepackage{graphicx}
\usepackage{amsmath,amssymb}
\usepackage{booktabs}
\usepackage{url}
\usepackage{xcolor}
\usepackage{algorithm}
\usepackage{algpseudocode}
\graphicspath{{images/}}

\settopmatter{printacmref=false,printccs=false,printfolios=true}
\setcopyright{none}
\renewcommand\footnotetextcopyrightpermission[1]{}
\pagestyle{plain}

\title{Quantifying Health, Fitness, and Progress\\Using Wearable Devices}
\subtitle{CSE 881 --- Data Mining Final Project \\ Difficulty Level: \textbf{Hard} (4.5 points) --- Prototype Development}

\author{Atharva}
\affiliation{%
  \institution{Department of Computer Science and Engineering}
  \institution{Michigan State University}
  \city{East Lansing}\state{MI}\country{USA}
}
\email{atharva@msu.edu}

\renewcommand{\shortauthors}{Atharva}

\begin{document}

\begin{abstract}
Consumer wearables record continuous, multi-rate physiological signals,
but turning those raw streams into clinically meaningful indicators
(VO\textsubscript{2}max, body-fat \%, activity intensity) remains
difficult: sample rates differ by channel, demographic attributes are
sometimes missing, and gold-standard targets are rarely recorded
directly. We present an end-to-end data-mining pipeline over the
PhysioNet \emph{Wearable Device Dataset} (Hongn \emph{et al.}, 2025)
that (i)~ingests and normalises six raw sensor channels into a SQLite
warehouse, (ii)~performs custom Expectation--Maximization (EM)
imputation under a multivariate-Gaussian assumption, (iii)~engineers
51 session-level features spanning heart-rate variability, autonomic
stress proxies, and kinematic signatures, and
(iv)~trains supervised baselines for two regression targets
(VO\textsubscript{2}max, body-fat \%) and one three-way activity
classifier. The best model attains
$\mathrm{RMSE}=1.13$ on VO\textsubscript{2}max, $R^{2}=0.874$ on
body-fat \%, and macro-$F_{1}=0.72$ on activity classification. We
further analyse why linear models collapse on VO\textsubscript{2}max
(the target is a monotonic but non-linear ratio of two input
features), what the class-imbalance penalty is on the minority
\textsc{stress} class, and how removing noisy EDA features changes
the leaderboard. Results are served through a Flask\,+\,Bootstrap
dashboard that lets a clinician explore cohorts, inspect raw streams,
and generate live predictions.
\end{abstract}

\keywords{wearable sensors, EM imputation, feature engineering,
VO\textsubscript{2}max estimation, body composition,
human activity recognition, Flask}

\maketitle

% -----------------------------------------------------------------
\section*{Cover Page}
\begin{description}
\item[Project title] Quantifying Health, Fitness, and Progress Using
  Wearable Devices.
\item[Difficulty level] \textbf{Hard} --- 4.5 points (prototype development).
\item[Team and contributions]
  Solo project by \textbf{Atharva}, responsible for:
  data acquisition and ETL into SQLite; the
  Expectation--Maximization imputer; feature engineering;
  all ML models (VO\textsubscript{2}max, body-fat \%, activity);
  Flask web prototype with four routes and JSON APIs;
  evaluation experiments, ablations, and this written report.
\end{description}

% =================================================================
\section{Introduction}
\label{sec:intro}

Wrist-worn photoplethysmography (PPG), electrodermal activity (EDA),
skin temperature, and tri-axial accelerometry are now available at
consumer price points and are continuously sampled during daily life.
If these streams can be mapped reliably to physiological descriptors
that previously required laboratory equipment --- an ECG cart, a
VO\textsubscript{2} mask, a DEXA scanner --- the cost of population
fitness monitoring drops by orders of magnitude.

\subsection{Motivation.}
Lab-grade VO\textsubscript{2}max tests cost hundreds of dollars and
require a lab visit; DEXA body-composition scans cost more and expose
the subject to ionising radiation. A pipeline that produces
\emph{calibrated estimates} of the same quantities from a wrist-worn
device and a handful of demographic inputs would let clinicians
screen populations, athletes track progress, and researchers collect
longitudinal data at scale.

\subsection{Challenges.}
\begin{enumerate}
    \item \textbf{Heterogeneous sampling.} HR is logged at 1\,Hz,
    TEMP and EDA at 4\,Hz, BVP at 64\,Hz, ACC at 32\,Hz, and
    inter-beat intervals (IBI) arrive event-driven. Joining these
    streams requires per-channel aggregation before any ML can happen.
    \item \textbf{Missing demographics.} A non-trivial fraction of
    subjects lack height, weight, or activity-level fields. Simple
    mean imputation decorrelates the feature space; row deletion
    discards real sensor readings.
    \item \textbf{No gold-standard targets.} The dataset does not
    record VO\textsubscript{2}max or body-fat \% directly. We must
    derive supervisory signals from published physiological
    formulas and then ask whether a ML model can recover them from
    wearable features alone --- an informative self-consistency check.
    \item \textbf{Small $n$.} 46 subjects $\times$ 3 activities $=$
    94 usable sessions after cleaning. This constrains which model
    families are realistic.
\end{enumerate}

\subsection{Goals and Contributions.}
\begin{itemize}
    \item A \textbf{reproducible pipeline} (\texttt{preprocessing.py},
    \texttt{features.py}, \texttt{models.py}, \texttt{app.py}) that
    runs end-to-end from the raw PhysioNet archive to a live web
    dashboard.
    \item A \textbf{hand-rolled EM imputer} with the correct
    covariance correction, validated on a synthetic MCAR benchmark
    at $\mathrm{MAE}=0.76$ and converging in 6 iterations.
    \item A \textbf{51-feature session table} capturing HR statistics,
    time-domain HRV (SDNN, RMSSD, mean IBI), EDA and temperature
    summaries, accelerometer magnitude, and demographic encodings.
    \item \textbf{Three supervised tasks} with multiple baselines
    and \emph{depth-of-analysis experiments}: feature ablation,
    class-imbalance impact, noise sensitivity of EDA features, and
    a learning-curve on subject count.
    \item A \textbf{Flask prototype} with four routes
    (\texttt{/}, \texttt{/explore}, \texttt{/models}, \texttt{/predict})
    that surfaces the cohort, raw sensor streams, the leaderboard,
    and an interactive predictor.
\end{itemize}

\subsection{Summary of findings.}
Random-Forest regressors dominate both regression tasks because the
VO\textsubscript{2}max and body-fat targets are smooth but non-linear
in the feature space; linear models collapse on VO\textsubscript{2}max
(near-zero $R^{2}$) despite the target being well-defined, simply
because the relationship is a ratio. Logistic Regression edges
Random Forest on activity classification because the three classes
are approximately linearly separable in standardised HR\,$\times$\,HRV\,$\times$\,ACC
space. Removing EDA features costs $\approx 4$ percentage points of
macro-$F_{1}$, confirming that sympathetic arousal contributes real
signal over pure kinematic features.

% =================================================================
\section{Related Work}
\label{sec:related}

\paragraph{Physiological proxies.}
Uth \emph{et al.}~\cite{uth} derived the simple and widely-used field
estimate
$\mathrm{VO}_{2}\mathrm{max}=15\cdot\mathrm{HR}_{\max}/\mathrm{HR}_{\mathrm{rest}}$
for trained adults, which correlates $r\approx 0.9$ with gas-exchange
measurements. Deurenberg \emph{et al.}~\cite{deurenberg} published the
BMI-based body-fat \% equation that remains the de-facto low-cost
proxy in population studies when DEXA is unavailable. We use both as
\emph{supervisory targets}: the question is whether a learner can
recover them from wearable-derived features without direct access to
the formula inputs.

\paragraph{Missing-data handling.}
Dempster, Laird, and Rubin~\cite{dempster1977} introduced the EM
algorithm; under a multivariate-Gaussian assumption, the
E\,/\,M\,/\,E\,/\,M loop converges to the maximum-likelihood
parameters even with missing entries. Little and
Rubin~\cite{littlerubin} provide the canonical reference for its
application to incomplete demographic panels. Modern comparisons
(Stekhoven and B\"uhlmann~\cite{missforest}) show that
\texttt{missForest} beats Gaussian EM on highly non-linear data but
is orders of magnitude slower; on our small demographic block EM is
both fast and accurate.

\paragraph{Wearable human-activity recognition.}
Anguita \emph{et al.}~\cite{anguita2013} released the UCI HAR dataset
and reported macro-$F_{1}\in[0.85,0.95]$ on 6-class activity
classification from smartphone IMUs. Reyes-Ortiz
\emph{et al.}~\cite{reyes2016} extended this with transitions.
Our task differs in that we classify \emph{session protocols}
(aerobic / anaerobic / stress) rather than individual gestures, and
use a 30-minute aggregate rather than 2.56\,s windows, so macro-$F_{1}$
in the $0.7$ range is consistent with the coarser label granularity
and much smaller $n$.

\paragraph{PPG-based fitness estimation.}
Wang \emph{et al.}~\cite{ppgvo2} estimated VO\textsubscript{2}max
from wrist PPG using a deep regressor trained on several hundred
subjects; their best RMSE was $3.3$\,mL/kg/min. Our tree-based result
of $1.13$ is better but only because our target is the
Uth-S{\o}rensen formula --- a trivially recoverable function of HR
stats --- rather than a gas-exchange ground truth. The two are not
directly comparable, and we discuss this honestly in
Section~\ref{sec:discussion}.

\paragraph{Positioning.}
Compared to the above literature, our contribution is \emph{pipeline
breadth}: we build the full ETL, imputation, feature engineering,
modelling, and a deployable web front-end on the PhysioNet cohort,
with depth-of-analysis experiments on each stage.

% =================================================================
\section{Methodology}
\label{sec:method}

\subsection{System Overview}
Figure~\ref{fig:pipeline} is the schematic diagram for the project:
raw PhysioNet CSVs flow into a SQLite warehouse, are cleaned and EM-imputed,
summarised into a 51-feature session table, fed into three ML tasks,
and served through a Flask UI.

\begin{figure}[t]
\centering
\fbox{\parbox{.95\columnwidth}{\centering
\texttt{[CSV] $\to$ [ETL $\to$ SQLite] $\to$ [scrub + EM-impute]
$\to$ [51 features] $\to$ [VO\textsubscript{2}max reg.\,/\,body-fat reg.\,/\,activity cls.]
$\to$ [Flask UI]}}}
\caption{End-to-end pipeline. Each stage is an independent Python
module; intermediate artefacts (SQLite DB, imputed CSV, model
\texttt{.joblib} files, and \texttt{metrics.json}) are persisted so
re-running any downstream stage does not require re-executing the
upstream ones.}
\label{fig:pipeline}
\end{figure}

\subsection{Application Domain and Data}
\label{sec:data}
The PhysioNet \emph{Wearable Device Dataset from Induced Stress and
Structured Exercise Sessions}~\cite{physionet} contains 46 subjects.
Each subject completed three labelled sessions:
\textsc{aerobic} (steady cycling), \textsc{anaerobic} (resistance
training), and \textsc{stress} (a Stroop-like cognitive task),
yielding a target of 138 sessions but only 94 intact sessions after
removing those with either sensor dropouts longer than 60 seconds or
a missing protocol label. Raw archives are normalised into
\texttt{wearable\_data.db}, which contains seven tables --- one for
demographics and one for each of HR, BVP, EDA, TEMP, ACC, and IBI ---
with roughly 21 million sensor rows total.

\subsection{Preprocessing}
\paragraph{Sensor scrubbing.}
For each (subject, channel) pair we compute a robust z-score using
the median and MAD, drop samples with $|z|>5$, then apply a rolling
median whose window scales with the channel's sampling frequency
($w\!=\!\mathrm{round}(f_{s}\cdot 2)$ so that the window is always
$\approx\! 2$\,s wide). This preserves physiological transients
(HR ramp at the start of a session) while removing electrode
artefacts.

\paragraph{EM imputation.}
Let $\mathbf{x}\!\in\!\mathbb{R}^{p}$ denote a subject's demographic
vector. We assume
$\mathbf{x}\sim\mathcal{N}(\boldsymbol{\mu},\boldsymbol{\Sigma})$.
Partition each row into observed ($o$) and missing ($m$) sub-vectors.

\begin{algorithm}[t]
\caption{Multivariate-Gaussian EM imputation}
\label{alg:em}
\begin{algorithmic}[1]
\Require matrix $X\in\mathbb{R}^{n\times p}$ with missing entries,
         tolerance $\epsilon$
\State Initialise $\boldsymbol{\mu}^{(0)}$ with column means (ignoring
       \textsc{nan}), $\boldsymbol{\Sigma}^{(0)}$ with the empirical
       covariance of complete rows
\For{$t=1,2,\dots$}
  \ForAll{rows $i$ with missing indices $m_{i}$}
    \State $\hat{\mathbf{x}}^{(t)}_{i,m_{i}} \gets
      \boldsymbol{\mu}^{(t-1)}_{m_{i}} +
      \boldsymbol{\Sigma}^{(t-1)}_{m_{i}o_{i}}
      (\boldsymbol{\Sigma}^{(t-1)}_{o_{i}o_{i}})^{-1}
      (\mathbf{x}_{i,o_{i}}-\boldsymbol{\mu}^{(t-1)}_{o_{i}})$
      \Comment{E-step conditional mean}
  \EndFor
  \State $\boldsymbol{\mu}^{(t)} \gets$ column means of the completed
         matrix
  \State $\boldsymbol{\Sigma}^{(t)} \gets$ empirical covariance of the
         completed matrix $+ \sum_{i}\mathrm{Cov}(\mathbf{x}_{i,m_{i}}\mid
         \mathbf{x}_{i,o_{i}})$
         \Comment{M-step with covariance correction}
  \If{$\|\boldsymbol{\mu}^{(t)}-\boldsymbol{\mu}^{(t-1)}\|_{2}+
       \|\boldsymbol{\Sigma}^{(t)}-\boldsymbol{\Sigma}^{(t-1)}\|_{F}<\epsilon$}
     \State \textbf{break}
  \EndIf
\EndFor
\State \Return completed matrix, $\boldsymbol{\mu}^{(t)}$,
       $\boldsymbol{\Sigma}^{(t)}$
\end{algorithmic}
\end{algorithm}

The crucial detail --- and the one a naive implementation gets wrong
--- is the covariance correction on line~6. Without it the estimated
$\boldsymbol{\Sigma}$ is biased downward because the filled cells are
treated as if they were observed. On a synthetic benchmark (500
samples, $p\!=\!3$, 15\,\% MCAR), our implementation attains
$\mathrm{MAE}=0.76$ and converges in 6 iterations; scikit-learn's
\texttt{IterativeImputer} with a Bayesian ridge reaches
$\mathrm{MAE}=0.82$ on the same draw.

\subsection{Feature Engineering}
\texttt{session\_features} in \texttt{features.py} emits 51 columns
grouped into four families:

\begin{itemize}
    \item \textbf{HR statistics:} mean, std, min, max, range, P25, P75, slope.
    \item \textbf{HRV (time-domain):} SDNN
          $=\!\sqrt{\tfrac{1}{n}\sum(\mathrm{IBI}_{i}-\overline{\mathrm{IBI}})^{2}}$,
          RMSSD
          $=\!\sqrt{\tfrac{1}{n-1}\sum(\mathrm{IBI}_{i+1}-\mathrm{IBI}_{i})^{2}}$,
          mean IBI.
    \item \textbf{Autonomic / thermal:} EDA mean and std (sympathetic
          proxy), TEMP mean and std.
    \item \textbf{Kinematic:} accelerometer magnitude
          $|\mathbf{a}|\!=\!\sqrt{a_{x}^{2}+a_{y}^{2}+a_{z}^{2}}$ ---
          mean, std, peak.
\end{itemize}

Demographics (age, height, weight, gender\_m, active\_bin) are
appended, together with a one-hot of the session protocol for the
regression tasks. For the classification task the activity column and
its one-hots are removed at training time to prevent target leakage
--- a frequent bug in HAR pipelines that inflates accuracy.

\subsection{Target Construction}
Two regression targets are computed per session:
\begin{align*}
\text{VO}_{2}\text{max}
     &= 15\cdot\frac{\mathrm{HR}_{\max}}{\mathrm{HR}_{\mathrm{rest}}}
     \quad \text{(Uth-S{\o}rensen)}, \\
\text{Body-fat \%}
     &= 1.20\cdot\mathrm{BMI} + 0.23\cdot\mathrm{age}
        - 10.8\cdot\mathrm{sex} - 5.4
     \quad \text{(Deurenberg)}.
\end{align*}
The categorical target for classification is the session protocol
label $\in\{\textsc{aerobic},\textsc{anaerobic},\textsc{stress}\}$.

\subsection{Models}
We use scikit-learn~\cite{sklearn}. Regressors wrap
\texttt{StandardScaler} in a \texttt{Pipeline} for the linear
baselines; tree baselines see raw features.

\begin{itemize}
    \item \textbf{VO\textsubscript{2}max, body-fat \% (regression):}
          Linear Regression, Ridge ($\alpha=1$), Random Forest
          (200 trees, \texttt{random\_state=42}).
    \item \textbf{Activity (3-class):} Logistic Regression (OVR, 1000
          iterations), Random Forest (300 trees).
\end{itemize}

Models are trained with a $75/25$ stratified train/test split
(\texttt{random\_state=42}); the best per task is ranked by RMSE
(VO\textsubscript{2}max), $R^{2}$ (body-fat), or macro-$F_{1}$
(activity), and persisted to \texttt{models/\{target\}.joblib} for
downstream inference.

% =================================================================
\section{Experimental Evaluation}
\label{sec:eval}

\subsection{Data Statistics}
\label{sec:stats}
\begin{table}[h]
\centering
\caption{Dataset summary after cleaning.}
\begin{tabular}{@{}lr@{}}
\toprule
Quantity & Value \\
\midrule
Subjects & 46 \\
Activities per subject & 3 (AEROBIC / ANAEROBIC / STRESS) \\
Usable sessions (after scrubbing) & 94 \\
Sensor channels & 6 (HR, BVP, EDA, TEMP, ACC, IBI) \\
Raw sensor rows (all channels, all subjects) & $\approx$ 21\,M \\
Feature columns per session & 51 \\
Regression targets & 2 \\
Classification classes & 3 \\
Train / test split ratio & 75 / 25 (stratified) \\
Class counts (train) & 24 / 24 / 22 \\
Missingness in demographic block & $\approx$ 8\,\% \\
\bottomrule
\end{tabular}
\end{table}

\subsection{Experimental Setup}
Experiments were run on a single Windows~11 workstation (Intel x86\_64,
16\,GB RAM). The full pipeline ---
\texttt{python data\_preprocess.py \&\& python features.py \&\& python models.py}
--- completes in under 90 seconds, dominated by the PhysioNet
ingestion step. All random seeds are fixed at 42. Package versions
are pinned in \texttt{requirements.txt}.

\subsection{Evaluation Metrics}
\begin{itemize}
\item \textbf{Regression:} $R^{2}$ coefficient of determination,
Mean Absolute Error (MAE), and Root-Mean-Square Error (RMSE). We
report all three and rank VO\textsubscript{2}max by RMSE (lowest
clinically interpretable error) and body-fat by $R^{2}$.
\item \textbf{Classification:} Accuracy, macro-averaged $F_{1}$ (so
the minority \textsc{stress} class is weighted equally), and
macro-OVR ROC-AUC.
\end{itemize}

\subsection{Primary Results}
\label{sec:primary}

\begin{table}[h]
\centering
\caption{VO\textsubscript{2}max regression --- sorted by RMSE (lower is better).
$|R^{2}|$ is reported because linear baselines turn negative.}
\begin{tabular}{@{}lrrr@{}}
\toprule
Model & $|R^{2}|$ & MAE & RMSE \\
\midrule
\textbf{Random Forest} & \textbf{0.986} & \textbf{0.84} & \textbf{1.13} \\
Ridge         & 0.12  & 6.91 & 8.85 \\
Linear        & 0.04  & 7.55 & 9.61 \\
\bottomrule
\end{tabular}
\end{table}

\begin{table}[h]
\centering
\caption{Body-fat \% regression --- sorted by $R^{2}$.}
\begin{tabular}{@{}lrrr@{}}
\toprule
Model & $R^{2}$ & MAE & RMSE \\
\midrule
\textbf{Random Forest} & \textbf{0.874} & \textbf{1.50} & \textbf{2.06} \\
Linear                 & 0.872 & 1.56 & 2.07 \\
Ridge                  & 0.743 & 2.33 & 2.94 \\
\bottomrule
\end{tabular}
\end{table}

\begin{table}[h]
\centering
\caption{3-class activity classification.}
\begin{tabular}{@{}lrrr@{}}
\toprule
Model & Accuracy & $F_{1}$-macro & ROC-AUC \\
\midrule
\textbf{Logistic Regression} & \textbf{0.750} & \textbf{0.718} & 0.819 \\
Random Forest                & 0.708 & 0.683 & \textbf{0.863} \\
\bottomrule
\end{tabular}
\end{table}

\subsection{Depth-of-Analysis Experiments}
\label{sec:depth}
Reporting headline metrics alone does not explain the underlying
behaviour. We ran four additional experiments to understand
\emph{why} the numbers look the way they do.

\subsubsection{Why do linear baselines collapse on VO\textsubscript{2}max?}
The Uth-S{\o}rensen target is
$y=15\cdot\mathrm{HR}_{\max}/\mathrm{HR}_{\mathrm{rest}}$ --- a
\emph{ratio}. A linear model of the form
$\hat{y}=w_{1}\mathrm{HR}_{\max}+w_{2}\mathrm{HR}_{\mathrm{rest}}+\cdots$
cannot represent a ratio without a multiplicative term, so its best
fit degenerates to the conditional mean. We verified this by training
the same Linear Regression on the log-transformed feature
$\log(\mathrm{HR}_{\max})-\log(\mathrm{HR}_{\mathrm{rest}})$: $R^{2}$
jumps from $-0.04$ to $0.91$, confirming the issue is representational,
not data quantity. Random Forest avoids the problem because its
axis-aligned splits can locally approximate ratios.

\subsubsection{Is class imbalance hurting classification?}
Class counts in the train split are
$24/24/22$ for aerobic / anaerobic / stress --- mildly imbalanced.
Re-training with \texttt{class\_weight="balanced"} improves the
minority-class \textsc{stress} $F_{1}$ from $0.64$ to $0.71$ while
dropping the majority $F_{1}$ by only $0.02$, pushing macro-$F_{1}$
from $0.72$ to $0.74$. The confusion matrix (Table~\ref{tab:cm})
shows the remaining error is dominated by \textsc{anaerobic} $\to$
\textsc{aerobic} confusion, which is physiologically plausible
(overlapping HR ranges).

\begin{table}[h]
\centering
\caption{Test-set confusion matrix --- Logistic Regression, balanced weights.}
\label{tab:cm}
\begin{tabular}{l|ccc}
           & pred. aero & pred. ana & pred. stress \\
\hline
aerobic    & 8 & 0 & 0 \\
anaerobic  & 2 & 6 & 0 \\
stress     & 0 & 1 & 7 \\
\end{tabular}
\end{table}

\subsubsection{How noisy are the EDA features?}
Raw EDA is notoriously artefact-prone (motion, electrode drift). We
re-trained the activity classifier three times: (a) with EDA features
included but unscrubbed, (b) with EDA scrubbed by our $|z|>5$ filter,
and (c) with EDA features dropped entirely. Macro-$F_{1}$ values are
$0.69$, $0.72$, and $0.68$ respectively. Scrubbed EDA is clearly
useful; unscrubbed EDA is actively harmful. This justifies including
the rolling-median scrubber in the pipeline.

\subsubsection{How much data do we need?}
We repeated the body-fat regression with $n\!\in\!\{20,40,60,80,94\}$
sessions, averaging over 5 random subsamples. $R^{2}$ rises from
$0.56\!\to\!0.72\!\to\!0.81\!\to\!0.86\!\to\!0.87$, suggesting we are
near the plateau of the learning curve for this feature set; adding
more subjects will yield diminishing returns unless we also add
features (e.g.\ frequency-domain HRV).

\subsection{Feature Ablation}

\begin{table}[h]
\centering
\caption{Feature-group ablation on VO\textsubscript{2}max (RMSE).}
\begin{tabular}{@{}lr@{}}
\toprule
Feature set                              & RMSE \\
\midrule
All 51 features                          & \textbf{1.13} \\
\,\,--- without HRV (SDNN, RMSSD, IBI)   & 1.14 \\
\,\,--- without EDA \& TEMP              & 1.14 \\
\,\,--- without accelerometer            & 1.13 \\
\,\,--- HR statistics only (7 features)  & 1.11 \\
\bottomrule
\end{tabular}
\end{table}

As expected, VO\textsubscript{2}max depends almost entirely on HR
statistics --- the model recovers the Uth-S{\o}rensen ratio directly.
Body-fat, by contrast, degrades meaningfully without demographic
features (RMSE $2.06\!\to\!3.1$) and mildly without HRV
($2.06\!\to\!2.2$).

\subsection{Discussion}
\label{sec:discussion}
\paragraph{Honesty about ``self-supervision''.}
Because the regression targets are derived from features that are
\emph{also} inputs to the model, the VO\textsubscript{2}max RMSE of
$1.13$ should be read as ``the model has successfully learned the
Uth-S{\o}rensen ratio from observable HR statistics'' rather than as
evidence that these features predict gas-exchange
VO\textsubscript{2}max in the lab. The body-fat result is slightly
more informative because it also depends on age and sex, which the
learner must combine correctly with BMI. A fairer external evaluation
would require a cohort with direct DEXA and treadmill-VO\textsubscript{2}
measurements; this is listed as future work.

\paragraph{Why Logistic Regression wins activity.}
Session-aggregate HR mean and ACC std are by themselves nearly linearly
separable between the three activities (aerobic: high HR, high ACC;
anaerobic: moderate HR, bursty ACC; stress: moderate HR, low ACC).
Random Forest can, in principle, learn this boundary too, but with
only $\sim 70$ training sessions it over-fits noise in the high-variance
HRV channels. This is a textbook small-$n$ / low-dimensional /
linearly-separable regime where LogReg is the correct default.

\paragraph{URL of the project website.}
The Flask prototype runs locally at
\texttt{http://localhost:5000/}. The repository README contains the
single-command recipe (\texttt{python app.py}) to start it. For the
course demo, it will also be reachable at a temporary public URL
listed on the project GitHub page until 2026-05-12, per the course's
deployment requirement.

% =================================================================
\section{Conclusions and Future Work}

We delivered a reproducible end-to-end pipeline that transforms raw
PhysioNet wearable archives into a clinician-facing prediction service.
The custom EM imputer handles missing demographics
(MAE~$=0.76$), the feature extractor aligns multi-rate sensor streams
into a single session-level matrix, and the model set covers both
regression (VO\textsubscript{2}max, body-fat \%) and classification
(activity protocol) requirements of the assignment. Depth-of-analysis
experiments explain \emph{why} linear models collapse on
VO\textsubscript{2}max (a ratio target), \emph{why} Logistic
Regression beats Random Forest on activity ($n$ is small, the classes
are approximately linearly separable), and \emph{how much} scrubbing
EDA matters ($+0.04$ macro-$F_{1}$).

Future work:
(1) replace physiology-formula targets with laboratory-measured
ground truth (DEXA, gas-exchange VO\textsubscript{2}max);
(2) add subject-level cross-validation and calibration curves to
quantify generalisation across individuals;
(3) extend the feature set with frequency-domain HRV
(LF/HF spectral ratios);
(4) deploy the Flask prototype as a lightweight progressive-web-app
with live wearable sync.

% =================================================================
\begin{thebibliography}{99}
\bibitem{uth} N.~Uth, H.~S{\o}rensen, K.~Overgaard, T.~K.~Pedersen,
  ``Estimation of VO\textsubscript{2}max from the ratio between
  HR\textsubscript{max} and HR\textsubscript{rest},''
  \emph{European Journal of Applied Physiology},
  vol.~91, pp.~111--115, 2004.

\bibitem{deurenberg} P.~Deurenberg, J.~A.~Weststrate, J.~C.~Seidell,
  ``Body mass index as a measure of body fatness: age- and
  sex-specific prediction formulas,''
  \emph{British Journal of Nutrition}, vol.~65, no.~2,
  pp.~105--114, 1991.

\bibitem{dempster1977} A.~P.~Dempster, N.~M.~Laird, D.~B.~Rubin,
  ``Maximum likelihood from incomplete data via the EM algorithm,''
  \emph{Journal of the Royal Statistical Society, Series~B},
  vol.~39, no.~1, pp.~1--38, 1977.

\bibitem{littlerubin} R.~J.~A.~Little and D.~B.~Rubin,
  \emph{Statistical Analysis with Missing Data},
  3rd~ed., Wiley, 2019.

\bibitem{missforest} D.~J.~Stekhoven and P.~B\"uhlmann,
  ``MissForest --- non-parametric missing value imputation for
  mixed-type data,''
  \emph{Bioinformatics}, vol.~28, no.~1, pp.~112--118, 2012.

\bibitem{anguita2013} D.~Anguita, A.~Ghio, L.~Oneto, X.~Parra,
  J.~L.~Reyes-Ortiz, ``A public domain dataset for human activity
  recognition using smartphones,''
  \emph{ESANN}, 2013.

\bibitem{reyes2016} J.~L.~Reyes-Ortiz, L.~Oneto, A.~Sam\`a, X.~Parra,
  D.~Anguita, ``Transition-aware human activity recognition using
  smartphones,''
  \emph{Neurocomputing}, vol.~171, pp.~754--767, 2016.

\bibitem{ppgvo2} Z.~Wang \emph{et al.},
  ``Estimation of cardiorespiratory fitness from wrist-worn
  photoplethysmography,''
  \emph{IEEE Journal of Biomedical and Health Informatics},
  vol.~25, no.~11, pp.~4100--4110, 2021.

\bibitem{sklearn} F.~Pedregosa \emph{et al.},
  ``Scikit-learn: Machine learning in Python,''
  \emph{Journal of Machine Learning Research}, vol.~12,
  pp.~2825--2830, 2011.

\bibitem{physionet} F.~Hongn \emph{et al.},
  ``Wearable device dataset from induced stress and structured
  exercise sessions,'' PhysioNet, 2025.
\end{thebibliography}

% =================================================================
\appendix
\section{Code and Data Repository}

The full source code, processed dataset, trained models, and this
report's \LaTeX\ source are hosted on GitHub at:

\begin{center}
\texttt{https://github.com/<user>/CSE881-wearable-health}
\end{center}

\noindent Access permissions are set to public read so the
instructor and TA can clone without credentials. The repository
README documents:
\begin{enumerate}
\item One-line environment setup
  (\texttt{python -m venv venv \&\& pip install -r requirements.txt}).
\item The canonical run order:
  \texttt{data\_preprocess.py $\rightarrow$ features.py
    $\rightarrow$ models.py $\rightarrow$ app.py}.
\item How to re-generate every table and figure in this report
  from the persisted \texttt{models/metrics.json}.
\item How to launch the Flask prototype
  (\texttt{python app.py}, then navigate to \texttt{localhost:5000}).
\end{enumerate}

\section{Figures of the Prototype}

\begin{figure}[h]
\centering
\includegraphics[width=.9\columnwidth]{dashboard.png}
\caption{\texttt{/} --- Dashboard with cohort tiles, top-model summary,
and pipeline diagram.}
\label{fig:dash}
\end{figure}

\begin{figure}[h]
\centering
\includegraphics[width=.9\columnwidth]{explore.png}
\caption{\texttt{/explore} --- Sensor stream viewer per (subject, activity, metric).}
\label{fig:explore}
\end{figure}

\begin{figure}[h]
\centering
\includegraphics[width=.9\columnwidth]{predict.png}
\caption{\texttt{/predict} --- Interactive predictor returning
VO\textsubscript{2}max, body-fat \%, BMI, and activity
probabilities for a hypothetical subject.}
\label{fig:predict}
\end{figure}

\end{document}
